\documentclass[12pt]{article}
\usepackage[utf8]{inputenc}
\usepackage[spanish]{babel}
\usepackage[a4paper, margin=2cm]{geometry}
\usepackage{tikz}
\usetikzlibrary{babel}
\usepackage[most]{tcolorbox}
\usepackage{amsmath}

% Usar fuente sin serifa para un aspecto más moderno (tipo infografía)
\renewcommand{\familydefault}{\sfdefault}

% Definición de colores con temática médica/científica
\definecolor{medblue}{RGB}{43, 108, 176}
\definecolor{medteal}{RGB}{49, 151, 149}
\definecolor{medlight}{RGB}{235, 248, 255}
\definecolor{meddark}{RGB}{26, 32, 44}
\definecolor{tearmuscle}{RGB}{220, 38, 38} % Rojo para el desgarro
\definecolor{restoring}{RGB}{16, 185, 129} % Verde para la fuerza restauradora

\begin{document}
\pagestyle{empty}

\begin{tcolorbox}[
    enhanced, 
    colback=medlight!30, 
    colframe=medblue, 
    title={\Large \textbf{Suma de Vectores: Estado de Equilibrio}}, 
    halign title=center, 
    fonttitle=\bfseries, 
    arc=4mm, 
    boxrule=1.5pt,
    shadow={2mm}{-2mm}{0mm}{black!30!white}
]

    \vspace{0.2cm}
    \begin{tcolorbox}[colback=white, colframe=medteal, arc=2mm, boxrule=1.2pt, title=\textbf{Caso Clínico / Problema}]
        Para neutralizar un desgarro muscular que tira con una fuerza de $\vec{F}_m = (25\hat{i} - 15\hat{j})$ N, un fisioterapeuta aplica una fuerza restauradora $\vec{F}_r = (-25\hat{i} + 15\hat{j})$ N. Demuestre matemáticamente el estado de equilibrio (Suma de fuerzas = 0).
    \end{tcolorbox}

    \vspace{0.4cm}
    
    \begin{center}
    \begin{tikzpicture}[scale=0.15]
        % Cuadrícula de fondo (escala de 5 en 5)
        \draw[step=5,gray!40,very thin] (-30,-20) grid (30,20);
        
        % Ejes cartesianos
        \draw[thick,->,meddark] (-32,0) -- (32,0) node[anchor=north west] {$x$ (N)};
        \draw[thick,->,meddark] (0,-22) -- (0,22) node[anchor=south east] {$y$ (N)};
        
        % Marcas numéricas
        \foreach \x in {-25,-15,-5,5,15,25} \draw (\x cm,1) -- (\x cm,-1) node[anchor=north, font=\scriptsize] {\x};
        \foreach \y in {-15,-5,5,15} \draw (1,\y cm) -- (-1,\y cm) node[anchor=east, font=\scriptsize] {\y};
        
        % Coordenadas principales
        \coordinate (O) at (0,0);
        \coordinate (Fm) at (25,-15);
        \coordinate (Fr) at (-25,15);
        
        % Vector 1: Fuerza Muscular (Desgarro)
        \draw[->, line width=2.5pt, tearmuscle] (O) -- (Fm) node[pos=0.8, below left, font=\bfseries, text=tearmuscle] {$\vec{F}_m$};
        \draw[dashed, tearmuscle, thick] (Fm) -- (25,0);
        \draw[dashed, tearmuscle, thick] (Fm) -- (0,-15);
        
        % Vector 2: Fuerza Restauradora (Fisioterapia)
        \draw[->, line width=2.5pt, restoring] (O) -- (Fr) node[pos=0.8, above right, font=\bfseries, text=restoring] {$\vec{F}_r$};
        \draw[dashed, restoring, thick] (Fr) -- (-25,0);
        \draw[dashed, restoring, thick] (Fr) -- (0,15);
        
        % Punto destacado (Origen / Punto de equilibrio)
        \filldraw[meddark] (O) circle (4pt) node[anchor=south west, font=\bfseries, fill=white, inner sep=1pt, xshift=2pt, yshift=2pt] {Lesión (0,0)};
        
    \end{tikzpicture}
    \end{center}

    \vspace{0.4cm}
    \begin{tcolorbox}[colback=medteal!10, colframe=medteal, arc=2mm, boxrule=1.2pt, title=\textbf{Desarrollo Matemático y Solución}]
        Para comprobar el estado de equilibrio, sumamos las componentes de todas las fuerzas involucradas. Si el sistema está en equilibrio, la fuerza total ($\vec{F}_{\text{total}}$) debe ser cero:
        \begin{align*}
            \vec{F}_{\text{total}} &= \vec{F}_m + \vec{F}_r = (25\hat{i} - 15\hat{j}) + (-25\hat{i} + 15\hat{j}) \\[1ex]
            \vec{F}_{\text{total}} &= (25 - 25)\hat{i} + (-15 + 15)\hat{j} = \mathbf{0\hat{i} + 0\hat{j} \text{ N}}
        \end{align*}
        \textbf{Resultado:} La suma de fuerzas es igual a \textbf{0 N}, por lo que el sistema se encuentra en un estado de equilibrio matemático y biomecánico.
    \end{tcolorbox}
    
\end{tcolorbox}

\end{document}